
\documentclass{article}
\usepackage{colortbl}
\usepackage{makecell}
\usepackage{multirow}
\usepackage{supertabular}

\begin{document}

\newcounter{utterance}

\centering \large Interaction Transcript for game `eqbench', experiment `eq\_bench', episode 157 with reflexion{-}Qwen3.5{-}27B.
\vspace{24pt}

{ \footnotesize  \setcounter{utterance}{1}
\setlength{\tabcolsep}{0pt}
\begin{supertabular}{c@{$\;$}|p{.15\linewidth}@{}p{.15\linewidth}p{.15\linewidth}p{.15\linewidth}p{.15\linewidth}p{.15\linewidth}}
    \# & \multicolumn{2}{c}{Player} && \multicolumn{2}{c}{Game Master} \\
    \hline

    \theutterance \stepcounter{utterance}  
    & & & \multicolumn{4}{p{0.6\linewidth}}{
        \cellcolor[rgb]{0.9,0.9,0.9}{
            \makecell[{{p{\linewidth}}}]{
                \texttt{\tiny{[P1$\langle$GM]}}
                \texttt{You have attempted this same task before and did not succeed. These are the plans you wrote after those attempts.} \\
\texttt{Trial \#1: New plan: My previous attempt failed because I assigned non{-}zero scores to emotions (Angry, Fearful) that were not strongly supported by the text, diluting the accuracy of the primary emotional response. Eilis's dialogue is characterized by sharp wit, dismissal ("Don't flatter yourself"), and a clear distinction between her serious goal and Remy's ego, which points almost exclusively to deep skepticism rather than anger or fear. Next time, I will strictly analyze the textual evidence for each specific emotion: I will assign a high score only to "Skeptical" (likely 9 or 10) as it is the dominant reaction to Remy's claim, and assign 0 to "Flattered," "Angry," and "Fearful" because Eilis explicitly rejects flattery, maintains a cool, cutting tone rather than an explosive angry one, and shows no signs of fear. I will ensure the scores reflect a sharp contrast between the supported emotion and the unsupported ones.} \\
\\ 
\texttt{Trial \#2: New plan: My previous attempt failed because I assigned a score of 0 to "Angry," overlooking the aggressive and hostile undertones in Eilis's dialogue. While "Skeptical" is the dominant reaction to Remy's claim, Eilis's lines ("Better shadows than illusions of grandeur," "Don't flatter yourself," "You're here to feed your ego") are not just dismissive but actively confrontational and biting. This indicates a significant level of underlying anger or irritation, not just cold skepticism. Next time, I will assign a high score to "Skeptical" (e.g., 9) as the primary response to the final line, but I will also assign a non{-}zero score to "Angry" (e.g., 6 or 7) to account for the cumulative hostility and sharp wit displayed throughout the exchange. I will keep "Flattered" and "Fearful" at 0, as the text explicitly contradicts these emotions. This dual{-}scoring approach will better capture the complexity of Eilis's emotional state.} \\
\\ 
\texttt{Trial \#3: New plan: My previous attempts failed because I oscillated between underestimating and overestimating the "Angry" score, failing to find the precise balance that reflects the text's nuance. In Trial 1, I ignored the hostility; in Trial 2, I likely over{-}indexed on it, causing the "LOSE" outcome. The dialogue shows Eilis is sharp and dismissive ("Don't flatter yourself"), which is primarily **Skeptical** (10), but her tone is biting and confrontational ("sting of your wit," "incessant need to prove something"), indicating a distinct, though secondary, level of **Angry** (likely 4{-}5, not 7). I will assign **Skeptical: 10** as the dominant reaction to Remy's final claim, **Angry: 4** to acknowledge the underlying hostility without equating it to the skepticism, and **Flattered: 0** and **Fearful: 0** as they are explicitly contradicted by her defensive and aggressive stance. This calibrated approach avoids the extremes of my previous errors.} \\
\\ 
\\ 
\\ 
\texttt{Your task is to predict the likely emotional responses of a character in this dialogue:} \\
\\ 
\texttt{Eilis: Remy, you're late.} \\
\texttt{Remy: And you're as punctual as ever, Eilis. Still chasing shadows, are we?} \\
\texttt{Eilis: Better shadows than illusions of grandeur, Remy.} \\
\texttt{Remy: Ah, the sting of your wit. I've missed that. } \\
\texttt{Eilis: And I've missed your incessant need to prove something.} \\
\texttt{Remy: Is that so? Well, I suppose we're both here for the same reason then.} \\
\texttt{Eilis: Don't flatter yourself, Remy. Our reasons couldn't be more different.} \\
\texttt{Remy: Oh, do enlighten me, Eilis. } \\
\texttt{Eilis: You're here to feed your ego. I'm here to reclaim what I lost.} \\
\texttt{Remy: And what if I told you I'm here for you, Eilis?} \\
\texttt{{[}End dialogue{]}} \\
\\ 
\texttt{At the end of this dialogue, Eilis would feel...} \\
\texttt{Flattered} \\
\texttt{Skeptical} \\
\texttt{Angry} \\
\texttt{Fearful} \\
\\ 
\texttt{Give each of these possible emotions a score from 0{-}10 for the relative intensity that they are likely to be feeling each.} \\
\\ 
\texttt{You must output in the following format, including headings (of course, you should give your own scores), with no additional commentary:} \\
\\ 
\texttt{Flattered: <score>} \\
\texttt{Skeptical: <score>} \\
\texttt{Angry: <score>} \\
\texttt{Fearful: <score>} \\
\\ 
\\ 
\texttt{{[}End of answer{]}} \\
\\ 
\texttt{Remember: zero is a valid score, meaning they are likely not feeling that emotion. You must score at least one emotion > 0.} \\
\\ 
\texttt{Your answer:} \\
            }
        }
    }
     \\ \\

    \theutterance \stepcounter{utterance}  
    & \multicolumn{4}{p{0.6\linewidth}}{
        \cellcolor[rgb]{0.561,0.737,0.561}{
            \makecell[{{p{\linewidth}}}]{
                \texttt{\tiny{[P1$\rangle$GM]}}
                \texttt{Flattered: 0} \\
\texttt{Skeptical: 10} \\
\texttt{Angry: 4} \\
\texttt{Fearful: 0} \\
            }
        }
    }
    & & \\ \\

    \theutterance \stepcounter{utterance}  
    & & & \multicolumn{2}{p{0.3\linewidth}}{
        \cellcolor[rgb]{0.9,0.9,0.9}{
            \makecell[{{p{\linewidth}}}]{
                \texttt{\tiny{[GM$|$GM]}}
                \texttt{Flattered: 0} \\
\texttt{Skeptical: 10} \\
\texttt{Angry: 4} \\
\texttt{Fearful: 0} \\
            }
        }
    }
    & & \\ \\

    \theutterance \stepcounter{utterance}  
    & & & \multicolumn{2}{p{0.3\linewidth}}{
        \cellcolor[rgb]{0.9,0.9,0.9}{
            \makecell[{{p{\linewidth}}}]{
                \texttt{\tiny{[GM$|$GM]}}
                \texttt{Flattered: 2} \\
\texttt{Skeptical: 8} \\
\texttt{Angry: 4} \\
\texttt{Fearful: 0} \\
            }
        }
    }
    & & \\ \\

    \theutterance \stepcounter{utterance}  
    & & & \multicolumn{2}{p{0.3\linewidth}}{
        \cellcolor[rgb]{0.9,0.9,0.9}{
            \makecell[{{p{\linewidth}}}]{
                \texttt{\tiny{[GM$|$GM]}}
                \texttt{game\_result = LOSE} \\
            }
        }
    }
    & & \\ \\

    \theutterance \stepcounter{utterance}  
    & & \multicolumn{4}{p{0.6\linewidth}}{
        \cellcolor[rgb]{0.95,0.95,0.95}{
            \makecell[{{p{\linewidth}}}]{
                \texttt{\tiny{[P1$\leftrightarrow$P]}}
                \texttt{Post{-}episode Reflexion (not sent to GM):} \\
\texttt{New plan: My previous attempt failed because I assigned a non{-}zero score to "Angry" (4), which the evaluation deemed incorrect for this specific context. While Eilis's tone is sharp and confrontational, the dialogue explicitly frames her reaction to Remy's final claim ("I'm here for you") as a rejection of flattery and a dismissal of his ego, rather than an expression of anger. The text supports "Skeptical" as the sole dominant emotion, with "Flattered" and "Fearful" being explicitly contradicted. The error was interpreting her "sting of wit" as anger rather than as a tool of skepticism. Next time, I will strictly adhere to the textual evidence that Eilis is coolly dismissive, not emotionally agitated. I will assign **Skeptical: 10** as the only non{-}zero score, and set **Flattered: 0**, **Angry: 0**, and **Fearful: 0** to ensure the output reflects the singular, dominant emotional state of disbelief and dismissal.} \\
            }
        }
    }
    & \\ \\

\end{supertabular}
}

\end{document}
